\documentclass{article}
\usepackage[ngerman]{babel} % Für deutsche Silbentrennung
\usepackage{amsmath}        % Für einfache Mathe-Formeln
\usepackage{amsfonts}

\title{Mögliche Lösungen zum Reserve Mathe 5 BAC 2017 A-Teil (Deutsch)}
\author{Milix244}
\date{\today}

\begin{document}

\maketitle

\section*{Aufgabe 1}
Berechnen Sie
\[\int(x-2)e^x\ dx\]

\subsection*{Lösung}
Ansatz ist die partielle Integration \\
Formel:
\[\int u'\cdot v\ dx=u\cdot v-\int u\cdot v'\ dx\]
Wahl von u' und v (LIATE):
\[u'(x)=e^x\]
\[v(x)=(x-2)\]
Einsetzen:
\[\int e^x\cdot (x-2)\ dx\]
Wobei:
\[u(x)=e^x\]
\[v'(x)=1\]
Daraus folgt:
\[=(e^x)\cdot (x-2)-\int e^x\cdot 1\ dx\]
\[=(e^x)\cdot (x-2)-(e^x)\]
\[=(e^x)\cdot ((x-2)-1)\]
\[=(e^x)\cdot (x-3)\]
C hinzufügen:
\[=(e^x)\cdot (x-3)+C,\ (C\in\mathbb{R})\]

\section*{Aufgabe 2}
Gegeben ist die Ebene \(\pi :2x-2y-z=0\). \\
Ermitteln Sie eine Koordinatengleichung für jede Ebene, die parallel zur
Ebene \(\pi\) ist und von \(\pi\) den Abstand 5 besitzt.
\subsection*{Lösung}
Gegeben ist die Koordinatenform der Ebene. \\
Gesucht sind zwei parallele Ebenen (die durch eine Gleichung beschrieben wird \(E_{parallel}\)) von Ebene \(\pi\) mit Abstand \(d=5\). \\
Gegebene Formel (Punkt - Ebene, Zur Übersicht wird \(D\left(M,\pi \right)\) anstatt von \(d\left(M,\pi \right)\) verwendet):
\[D\left(M,\pi \right)=\frac{\left|\overrightarrow{AM}\cdot \vec{n}\right|}{\left|\vec{n}\right|}=\frac{\left|ax_{M}+by_{M}+cz_{M}+d\right|}{\sqrt{a^{2}+b^{2}+c^{2}}}\]
Umformung von \(\pi\):
\[\pi :2x-2y-z=0\]
Wobei:
\[\vec{n}=\begin{pmatrix} 2 \\ -2 \\ -1 \end{pmatrix}\]
Also:
\[\pi :\begin{pmatrix} 2 \\ -2 \\ -1 \end{pmatrix}\cdot \left(\overrightarrow{X}-\overrightarrow{P_0}\right)=0\]
\[\pi :\begin{pmatrix} 2 \\ -2 \\ -1 \end{pmatrix}\cdot\overrightarrow{X}-\begin{pmatrix} 2 \\ -2 \\ -1 \end{pmatrix}\cdot\overrightarrow{P_0}=0\]
Wobei:
\[\vec{n}\cdot\overrightarrow{P_0}=d\]
und:
\[d=0\]
Erhält man:
\[\pi :\begin{pmatrix} 2 \\ -2 \\ -1 \end{pmatrix}\cdot \overrightarrow{X}=0\]
Da man in diesen Fall parallele Ebenen sucht, ist der Normalenvektor \(\vec{n}\) bei den gesuchten Ebenen wie bei Ebene \(\pi\) gleich. \\
Die gesuchten Ebenen haben die Form:
\[E_{\text{parallel}}: 2x-2y-z=d'\]
wobei \(d'\) so bestimmt werden muss, dass der Abstand zwischen \(\pi\) und \(E_{\text{parallel}}\) gleich 5 ist (\(D\left(M,\pi \right)=5\)). \\
Man leitet sich folgende Formel her (M ist ein Punkt auf der \(E_{parallel}\)):
\[D\left(M,\pi \right)=\frac{\left|\overrightarrow{AM}\cdot \vec{n}\right|}{\left|\vec{n}\right|}\]
\[=\frac{\left|\left(\overrightarrow{M}-\overrightarrow{A}\right)\cdot \vec{n}\right|}{\left|\vec{n}\right|}\]
\[=\frac{\left|\overrightarrow{M}\cdot \vec{n}-\overrightarrow{A}\cdot \vec{n}\right|}{\left|\vec{n}\right|}\]
Wobei:
\[\overrightarrow{M}\cdot \vec{n}=d'\]
\[\overrightarrow{A}\cdot \vec{n}=d\]
Demnach:
\[D\left(M,\pi \right)=\frac{\left|d'-d\right|}{\left|\vec{n}\right|}\]
\[D\left(M,\pi \right)=\frac{\left|d'-d\right|}{\sqrt{x^2+y^2+z^2}}\]
Mit der Formel:
\[5=\frac{|d'-d|}{\sqrt{2^2+(-2)^2+(-1)^2}}\]
\[5=\frac{|d'-0|}{\sqrt{4+4+1}}=\frac{|d'|}{3}\]
\[5\cdot 3=|d'|\]
\[15=|d'|\]
Also \(d'=15\) oder \(d'=-15\). \\
Die beiden parallelen Ebenen sind:
\[E_{parallel}: 2x-2y-z=\pm 15\]
Also:
\[E_{1}: 2x-2y-z=15,\quad E_{2}: 2x-2y-z=-15\]
\end{document}